\chapter{Implementation}\label{s:implementation}

This chapter describes how the system designed in Chapter~\ref{s:design} is
realised in code. Following the component overview of
Chapter~\ref{s:overview}, the pipeline is implemented as a Python
application with two entry points: a command-line interface and an
interactive dashboard. Both entry points share the same orchestrator,
retrieval stack, and semantic cache, while fine-tuning runs separately as a
batch job on Snellius. Across these components, run artefacts are organised
under a single \texttt{outputs/} directory with dedicated subdirectories for
models, evaluation results, training metrics, pipeline results, generated
BraneScript, and logs.

\section{Pipeline Orchestrator}\label{s:pipeline-impl}

The orchestrator (\texttt{pipeline.py}) loops over the
generation--validation--execution cycle described in
Section~\ref{s:overview} for at most three attempts per request. Before
attempting generation, it queries the semantic cache
(Section~\ref{s:cache-impl}) for the intent and returns the cached script and
result immediately on a hit; on a miss, it decomposes the intent into
sub-tasks via \texttt{IntentDecomposer} before retrieval, as described in
Section~\ref{s:decomposition}. This step reuses the generation model already loaded in
memory rather than loading a second one, so its cost is one extra generation
call rather than one extra model; a \texttt{--skip-decomposition} flag in the orchestrator
disables the step outright, isolating its cost from the rest of the pipeline so its
effect on the generated BraneScript can be assessed independently of retrieval and
generation. A candidate workflow is checked for valid function
and argument usage against the retrieved package context before
\texttt{execute\_workflow()} ever submits it to Brane, which compiles and
executes it as a single command and returns \texttt{stdout},
\texttt{stderr}, and an exit code that distinguishes a compilation failure
from a runtime failure. On success, the (intent, script, result) triple is
written back to the semantic cache and also persisted to disk for later
inspection and reuse. Intents the
orchestrator judges too ambiguous to act on are handled separately: rather
than attempting generation, it returns a clarifying question to the user.

Setting up a pipeline run once at start-up rather than per-request loads
the embedding model, builds the single package retriever shared across the
run, and connects to the semantic cache's ChromaDB collection; if that
connection fails, the failure is caught and logged and the run proceeds
with caching disabled rather than failing outright, while the BraneScript
syntax reference (\texttt{data/syntax\_reference.md}) is loaded once as plain text and needs no embedding model
at all: at roughly 450 lines (~2,000 tokens) it is small enough to be
included in full in every prompt without approaching context-length limits. A command-line flag allows the cache to be bypassed explicitly for
a given invocation.

\section{Knowledge Base and Retrieval}\label{s:kb-impl}

\texttt{src/knowledgeBase.py} builds and queries the single vector store
introduced in Section~\ref{s:retrieval-design}, called \texttt{brane\_pkg\_db}, a
\href{https://www.trychroma.com/}{ChromaDB} collection persisted to disk and
accessed through LangChain's \texttt{Chroma} and
\texttt{HuggingFaceEmbeddings} wrappers, the same LangChain abstractions used
throughout the retrieval stack. Package specifications
(defined as \texttt{container.yml} files following Brane package definition) 
are parsed and exploded into per-function
\texttt{Document} chunks, each annotated with the owning package name and a
short usage hint, and embedded with \texttt{all-MiniLM-L6-v2}. 
On top of that, each package's \texttt{keywords.txt} file (Section~\ref{s:brane-packages}) is
loaded as its own short, unchunked document tagged with the owning package
name; \texttt{PkgRetriever} reads these documents back out of the collection
at startup and builds an in-memory alias table from them, so the hybrid
retrieval strategy described in Chapter~\ref{s:design} never has to guess a
package's domain vocabulary from its function signatures alone. BraneScript syntax is
instead finally loaded via \texttt{data/syntax\_reference.md}, a curated file of BraneScript constructs
of interest to the generation model, which is small enough  (approximately 2,000 tokens) to be included in full and
formatted directly into the generation system prompt for every request.

As briefly mentioned in (Section~\ref{s:retrieval-design}), in order to support the 
continuous development of the systems, the knowledge base does not need
to be rebuilt from scratch when new package or dataset content is added.
\texttt{knowledgeBase.py} allows the addition of embeddings directly to the existing \texttt{brane\_pkg\_db}
collection, so registering a new package only requires embedding that
package's own documents and inserting them, leaving every previously
indexed package untouched. This means that adding a new
package, or editing an existing package's \texttt{container.yml},
\texttt{keywords.txt}, or \texttt{README.md}, only costs the time to embed
the changed content itself and not the whole catalogue. 
The current implementation of the pipeline doesn't detect a stale knowledge base automatically, so re-running the
update is still a manual part of the package-authoring workflow rather
than one triggered by pipeline startup.

\section{Brane Packages}\label{s:brane-packages}

The seven packages implemented under \texttt{packages/} as the evaluation
testbed for this thesis were purpose-built to provide a concrete end-to-end
domain on which to evaluate the system.

Every package follows the same structure: a \texttt{container.yml}
declaring the package's functions, their typed input parameters and return values,
and any custom \texttt{class} types it exposes to BraneScript; a single
Python source file implementing the actual logic; and a \texttt{keywords.txt}
listing the domain vocabulary a scientist might use to refer to the package
without naming it directly (Section~\ref{s:retrieval-design}). Additionally,
the packages include a \texttt{README.md} documenting, in
prose, the package's Brane dispatch conventions and the meaning of each
exposed function and class, written from a human maintainer with the goal of extending the package documentation.

Each package follows a common convention: the
container's entry point reads the invoked function name from
\texttt{sys.argv[1]} (declared in the package's \texttt{command.args}
specification), rather than relying on environment variables, and routes to
the corresponding Python function. Structured return values are serialised as
tagged class instances (\texttt{["ClassName", \{...\}]}), consistent with the
typed data representation decision in Chapter~\ref{s:design}. Every package is
of Brane's \texttt{ecu} (\emph{executable code unit}) kind, for which Brane
itself generates the Docker container image from the \texttt{dependencies},
\texttt{install}, and \texttt{files} fields declared in
\texttt{container.yml}.

\section{Fine-Tuning Infrastructure}\label{s:finetune-impl}

\subsection{Training the models}

\texttt{src/fine\_tuning/train.py} runs fine-tuning with the base model
supplied as a configuration value rather than hard-coded, so the same training code
produced the Qwen3.5-9B and Qwen3.6-27B adapters compared in
Chapter~\ref{s:evaluation} without any per-model changes. Fine-tuning was
limited to these two model sizes due to GPU memory and allocation constraints
on Snellius: training a 27B model with 4-bit quantisation and LoRA already
requires a full H100 GPU (80\,GB VRAM), leaving larger models (70B and above)
outside the available resource envelope for this project. Default
hyperparameters are LoRA rank~16 (alpha~32), 3 epochs, an effective batch
size of~8, and a learning rate of $2\times10^{-4}$ on a cosine schedule,
using 4-bit quantisation through Unsloth where available and a plain
Hugging Face/PEFT fallback otherwise; the response-only loss masking
described in Chapter~\ref{s:design} is applied at this stage. Training runs
on Snellius with a single H100 GPU and resumes automatically from the latest checkpoint
if a job is pre-empted, rather than restarting from the base model. 
Every run's full hyperparameter set and training progress are logged
alongside the resulting adapter so that multiple SFT runs against the same
base model remain distinguishable later.

\subsection{Tokenisation and Sequence Length}

Training examples are formatted using the same prompt template as inference:
a system prompt followed by the user intent, with the BraneScript workflow as
the expected assistant response. The maximum sequence length is set to 2048
tokens, which accommodates the full prompt (system prompt + syntax reference
+ retrieved package context + intent) plus the expected BraneScript output
without truncation for all but the longest multi-step workflows in the
training set. Truncation would silently corrupt training examples by cutting
off the BraneScript output mid-workflow (as experienced in some runs),
producing malformed targets that the model would learn to reproduce; the 2048-token limit was chosen by inspecting
the length distribution of the training set and confirming that all examples
fit within it.

Response-only loss masking is applied so that the model is penalised only on
the BraneScript output tokens and not on the prompt. Without this, the model
would spend training capacity learning to reproduce the (fixed) system prompt
and retrieved context, diluting the signal from the BraneScript generation
task itself.

\subsection{Adapter Merging}

After training, the LoRA adapter weights are merged into the base model
weights using \texttt{model.merge\_and\_unload()}, producing a single
standalone model checkpoint that can be loaded for inference without the PEFT
library. This avoids the runtime overhead of applying the adapter at each
forward pass and simplifies deployment: the merged model is identical in
interface to the base model and can be loaded with the standard Hugging Face
\texttt{AutoModelForCausalLM} API. The unmerged adapter is retained alongside
the merged checkpoint so that further fine-tuning runs can resume from the
adapter state rather than re-merging.

Evaluation is handled by \texttt{evaluate.py}, which computes the three
metrics used throughout Chapter~\ref{s:evaluation} and described in Subsection~\ref{s:eval-metrics}.
Each evaluation run records which training
configuration produced the model under test, so that the dashboard
(Section~\ref{s:frontend}) can distinguish between multiple SFT runs of
the same base model without additional bookkeeping.

\section{Semantic Cache}\label{s:cache-impl}

\subsection{Embedding Model}

Intent similarity is computed using \texttt{all-MiniLM-L6-v2}, a 22M-parameter
sentence-embedding model that maps text to 384-dimensional unit vectors.
Cosine similarity between the query and stored vectors is used as the
retrieval score. The model was chosen for its inference speed and semantic
quality on short English sentences, which matches the typical intent length
of 10--30 words. Its known limitation, which revolved around encoding
topic rather than precise parameter values, is the direct cause of the 
false-positive floor noticed in Section~\ref{s:eval-cache} and further discussed in Section~\ref{s:disc-rq4}.

\subsection{Persistence}

The cache collection is stored on disk and survives server restarts without
needing to be rebuilt. The 584 seeded intents and any workflows added during
interactive use are retained across process boundaries, meaning the first
request for a known intent always hits the cache rather than requiring a
prior successful generation to populate it.
The cache collection is thought to be deployed to remote servers, but the current implementation
runs it locally on the same machine as the dashboard and orchestrator, which is sufficient for
the single-user research deployment.

\subsection{Concurrency}

The cache does not implement explicit application-level locking. The
underlying storage engine serialises writes internally, so concurrent stores
will not corrupt the collection, but a lookup that begins just before a
concurrent store completes may miss the new entry. In the current deployment
this race window is narrow, as the only runtime caller of store is the
interactive entry point handling one request at a time. A multi-user
production deployment would require either an explicit lock or a migration to
a client-server storage mode.

\subsection{Cache Invalidation}

The cache exposes an invalidation operation that removes all stored entries
whose similarity to a supplied intent exceeds the current threshold. The
intended use case is package maintenance: if a new version of a Brane package
changes a function signature, any cached workflow calling the old signature
will silently produce a compile failure on the next execution. Invalidating
with a representative intent for the affected package forces fresh generation
for subsequent requests. This is a manual operation in the current
implementation; automatic invalidation on package version change is left as
future work.

The cache is seeded with all 584 training intents at startup, so the first
request for any known intent always hits the cache. The evaluation pipeline
used throughout Section~\ref{s:evaluation} deliberately bypasses the cache entirely, 
as its purpose is to measure raw model performance rather than cache behaviour.

\section{Dashboard}\label{s:frontend}

The dashboard backend is implemented with Flask. The frontend is a
single-page application served as static files by the same process, keeping
the deployment to a single \texttt{python server.py} invocation on the
researcher's machine.

Snellius job submission is handled over SSH via subprocess calls: the server
connects to Snellius using a key-based SSH configuration read from environment
variables and invokes the SLURM batch script remotely. The status polling works
the same way: the server queries the SLURM job queue over SSH at each poll
interval and maps the scheduler state to the dashboard's own
Queued~$\rightarrow$~Generating~$\rightarrow$~Executing~$\rightarrow$~Done
progression. No persistent connection to Snellius is maintained and each status
check is an independent SSH call.

Results are finally communicated back through a shared file-based job directory: the
generation worker on Snellius writes its output to a known path, and the
dashboard backend reads it on the next poll. This decoupling contributed to the design decision that the
dashboard does not need to stay connected to Snellius while the job runs.

\section{Job Watcher}\label{s:job-watcher}

The job watcher is a lightweight polling daemon that runs on whichever
machine has network access to the Brane instance. In the current
deployment, it is the researcher's own machine rather than the Snellius compute
node that ran generation, since Snellius does not support Docker. 
This separation is intentional: generation (the
LLM call and decomposition) happens on the HPC cluster, but Brane execution
must happen where Brane is installed and reachable, which need not be the
same machine.

The daemon monitors a shared job directory for BraneScript files deposited
by the generation pipeline. For each pending script it invokes the Brane
compiler and runtime directly, so the exit code and output are always Brane's
own verdict rather than a re-implementation of its compilation or execution
logic. The result, whether success or failure, standard output, and any error
messages, is written back to a result file that the dashboard picks up and,
on failure, feeds into the retry sequence shown in
Figure~\ref{fig:sequence-diagram}. Keeping execution separate from the
dashboard process limits the dashboard's responsibilities to submission,
polling, and presentation, and makes it possible to run Brane on a different
host from the server without changing any pipeline logic.

\section{Interaction Sequence}\label{s:interaction-sequence}

Figure~\ref{fig:pipeline-architecture} shows the components involved in a
single request but not the timing of the messages between them, in
particular the asynchronous submit/poll/retry pattern of the interactive
dashboard entry point. Figure~\ref{fig:sequence-diagram} makes that timing
explicit for the case a researcher exercises in day-to-day use: the
researcher's browser, the dashboard backend, the semantic cache, the
generation worker running as a batch job on Snellius, and the Brane runtime
that ultimately compiles and executes the workflow.

The diagram distinguishes two branches. On a \textbf{cache hit} (top frame,
steps 1--3), the dashboard looks up the submitted intent against the
semantic cache and, on a similarity match, returns the cached script and its
previously recorded execution result directly to the researcher; nothing on
Snellius or Brane is touched, and no model is loaded. On a \textbf{cache
miss} (bottom frame, steps 4--14), the dashboard instead submits a SLURM
batch job to Snellius and immediately hands the researcher back a request
identifier and a ``submitted'' status (steps 5--6), rather than blocking the
HTTP request for however long generation and execution take; the
researcher's browser polls a status endpoint for that identifier until the
job completes (step 7). On the Snellius side, the worker performs the
process-architecture steps described in Section~\ref{s:pipeline-variants} in
order -- decomposing the intent, retrieving package/dataset context, and
generating a candidate script (step 8) -- and then submits that script to
the Brane runtime with a single \texttt{brane workflow run} call (step 9),
receiving back an exit code and captured stdout/stderr (step 10). The worker
writes this result where the dashboard can find it (step 11); if execution
failed and retries remain, the dashboard resubmits the same request with the
failure fed back into the prompt as corrective context (step 12), returning
to step 5. Once a script executes successfully, the result is written to
the semantic cache (step 13) before the final script and outcome are
returned to the researcher (step 14).

The command-line entry point (\texttt{pipeline.py}) follows the same overall
shape -- cache lookup, decomposition, retrieval, generation, cache store on
success -- but differs from the diagram in one respect: it has no
equivalent of step~7's polling, since it runs synchronously to completion on
the researcher's own machine instead of via a remote batch scheduler.

\begin{figure}[htbp]
    \centering
    \resizebox{\textwidth}{!}{%
    \begin{tikzpicture}[
        lifeline/.style={draw, thick, minimum width=2.6cm, minimum height=0.9cm,
            align=center, font=\small\bfseries, fill=gray!12},
        msg/.style={-{Latex[length=2.2mm]}, thick},
        ret/.style={-{Latex[length=2.2mm]}, thick, dashed},
        fb/.style={-{Latex[length=2.2mm]}, thick, dashed, red!65!black},
        lbl/.style={font=\scriptsize, align=center},
        seqframe/.style={draw, thick, rounded corners=1pt, inner sep=6pt},
        frametitle/.style={draw, thick, fill=yellow!20, font=\bfseries\scriptsize,
            align=left, anchor=north west, inner sep=3pt}
    ]
        \node[lifeline] (usr)   at (0, 0)    {Researcher};
        \node[lifeline] (dash)  at (3.8, 0)  {Dashboard};
        \node[lifeline] (cache) at (7.6, 0)  {Semantic\\Cache};
        \node[lifeline] (gen)   at (11.4, 0) {Snellius\\worker};
        \node[lifeline] (brane) at (15.2, 0) {Brane\\runtime};

        \foreach \p in {usr, dash, cache, gen, brane}
            \draw[dotted] (\p.south) -- ++(0,-20.6);

        \draw[msg] (usr.south |- 0,-1.3) -- node[lbl, above]{1: submit intent} (dash.south |- 0,-1.3);

        \draw[msg] (dash.south |- 0,-2.9) -- node[lbl, above]{2: lookup(intent)} (cache.south |- 0,-2.9);
        \draw[ret] (cache.south |- 0,-3.7) -- node[lbl, above]{2': cached (script, result)} (dash.south |- 0,-3.7);
        \draw[msg] (dash.south |- 0,-4.9) -- node[lbl, above]{3: return cached script + result} (usr.south |- 0,-4.9);

        \node[seqframe, fit={(usr.south |- 0,-2.2) (cache.east |- 0,-5.3)}] (frameA) {};
        \node[frametitle] at (frameA.north west) {Cache hit (similarity $\geq 0.92$)};

        \draw[msg] (dash.south |- 0,-6.9) -- node[lbl, above]{4: lookup(intent) --- miss} (cache.south |- 0,-6.9);
        \draw[msg] (dash.south |- 0,-8.0) -- node[lbl, above]{5: submit SLURM job (sbatch)} (gen.south |- 0,-8.0);
        \draw[msg] (dash.south |- 0,-9.1) -- node[lbl, above]{6: req\_id, status=submitted} (usr.south |- 0,-9.1);
        \draw[msg] (usr.south |- 0,-10.2) -- node[lbl, above]{7: poll GET /api/generate/req\_id (repeated)} (dash.south |- 0,-10.2);

        \draw[msg] (gen.south |- 0,-11.3) -- ++(2.0,0) -- ++(0,-1.3) --
            node[lbl, right, text width=8cm]{8: decompose intent into sub-tasks; retrieve pkg/dataset context; generate BraneScript}
            (gen.south |- 0,-12.6);

        \draw[msg] (gen.south |- 0,-14.0) -- node[lbl, above]{9: brane workflow run} (brane.south |- 0,-14.0);
        \draw[ret] (brane.south |- 0,-14.9) -- node[lbl, above]{10: exit code, stdout/stderr} (gen.south |- 0,-14.9);
        \draw[msg] (gen.south |- 0,-16.0) -- node[lbl, above]{11: write result} (dash.south |- 0,-16.0);
        \draw[fb]  (dash.south |- 0,-17.5) -- node[lbl, above, text width=5.5cm]{12: on failure, attempts remain: resubmit with error feedback ($\to$ step 5)} (gen.south |- 0,-17.5);
        \draw[msg] (dash.south |- 0,-18.8) -- node[lbl, above]{13: on success: store(intent, script, result)} (cache.south |- 0,-18.8);
        \draw[msg] (dash.south |- 0,-19.9) -- node[lbl, above]{14: final script + execution outcome} (usr.south |- 0,-19.9);

        \node[seqframe, fit={(usr.south |- 0,-6.2) (gen.east |- 0,-20.3)}] (frameB) {};
        \node[frametitle] at (frameB.north west) {Cache miss --- generate, execute, retry ($\leq 3$ attempts)};
    \end{tikzpicture}}
    \caption[Interaction sequence for the interactive dashboard entry point]{Sequence of interactions for the interactive (dashboard) entry
    point, discussed in the text above.}
    \label{fig:sequence-diagram}
\end{figure}

\FloatBarrier

%%% Local Variables:
%%% mode: latex
%%% TeX-master: "../thesis"
%%% End:
